\bookchapter{الملحق E: شكر وتوضيحات}

شكرًا لإصرارك على قراءة هذا «الكورس الهاتفي» حتى نهايته! أثناء الكتابة، استعنّا بعدد كبير من الأدلة الرسمية للشركات المصنّعة وأدلة المستخدمين والمقالات التقنية التوعوية المتاحة للجمهور، لضمان دقة المعلومات.

التقنية تتطور يوماً بعد يوم، ومحتوى الكتاب مبني على التقنيات والمنتجات السائدة وقت التأليف. وعندما تقرؤه، ربما ظهرت ميزات ومنتجات جديدة. نأمل أن تساعدك الأفكار والأساليب التي يقدمها هذا الكتاب على التعامل بثقة مع المزيد من الأجهزة الرقمية الجديدة في المستقبل.

\booksection{أبرز مستجدات الإصدار 1.6}

\begin{itemize}
  \item توسيع «قيمة الأصول الرقمية وفلسفة إدارتها»، وإضافة فصل «إدارة الصور ونظام الملفات» (التسمية بالزمن+الحدث، والتصنيف والترتيب، والتنظيف التلقائي بعد 72 ساعة، وإخفاء بيانات EXIF الحساسة).
  \item استكمال منظومة الحماية بـ«الوضع النقي/وضع الحماية» و«الطفيليات الإلكترونية».
  \item إضافة أشكال الاحتيال الشائعة في عصر الذكاء الاصطناعي (تبديل الوجوه وتقليد الأصوات / انتحال صفة جهات رسمية / استثمارات وهمية) وأفكار الوقاية والخروج منها.
  \item إعادة كتابة الوقاية من هجمات «فتح الصندوق» (Doxxing) والهندسة الاجتماعية.
  \item توسيع فصل إمكانية الوصول (عمليات VoiceOver/TalkBack الكاملة، والتكبير وفلاتر الألوان، ووصف الصور بالذكاء الاصطناعي، والترجمة النصية الفورية، وقناتا المكالمات المزدوجتان RTT/TTY).
  \item إعادة صياغة إطار كتابة الأوامر التوضيحية بمنهج «الدور-الخلفية-المهمة-القيود»، مع 3 حالات كاملة من سيناريوهات واقعية.
  \item إضافة ثلاثة ملاحق تقنية متقدمة: ملحق الفصل 1: تحليل توسعة ذاكرة هواتف Android (التعريف الجوهري، مبدأ zRAM/swap، عنق زجاجة السرعة، حدود حالات الاستخدام، وتسمية كل شركة وطريقة التفعيل).
  \item ملحق الفصل 20: Termux — بيئة طرفية Linux محمولة على أجهزة Android (الوظائف الأساسية، منظومة الإضافات، التثبيت والإعداد، وسيناريوهات نموذجية مثل تطوير تطبيقات الهاتف وإدارة الخوادم والاختبار الأمني).
  \item ملحق الفصل 23: دليل المبتدئين إلى GitHub — من استضافة الشيفرة إلى منصة التعاون (التمييز بين Git وGitHub، المفاهيم الأساسية، التسجيل والبدء، سير التعاون عبر Pull Request، تنبيهات أمنية، وخيارات الوصول من داخل الصين).
\end{itemize}

\booksection{أبرز مستجدات الإصدار 1.7}

أولًا: إضافة 4 فصول جديدة

\begin{itemize}
  \item الفصل 29: الصحة الرقمية — تحليل وقت استخدام الشاشة، وحدود التطبيقات، ووضع عدم الإزعاج، وقاعدة 20-20-20 لحماية العينين، وإنقاذ الرقبة، وإدارة النوم، و«العمى الانتقائي» لقلق المعلومات، وتحدي أسبوع الصحة الرقمية.
  \item الفصل 30: مشاركة الشاشة والسيارة — مقارنة بين معسكرات AirPlay / Miracast / DLNA الثلاثة، وحلول التوصيل السلكي، ووضع سطح المكتب Samsung DeX، وحلول السيارة الأربعة CarPlay / Android Auto / HiCar / CarLife، وتطبيق عملي لمشاركة الشاشة في قاعات الاجتماعات.
  \item الفصل 31: صناعة الفيديوهات القصيرة — الأسئلة الثلاثة لاختيار الموضوع، وإعدادات التصوير في 5 خطوات، وخمس حركات كاميرا (الدفع والسحب والتأرجح والإزاحة والملاحقة)، وسير عمل CapCut في 6 خطوات، وأساسيات صناعة الغلاف، ومنهجية «أول 3 ثوانٍ تحدد المصير».
  \item الفصل 32: دليل الطوارئ الأمنية — إجراءات الساعة الثلاثين الذهبية عند فقدان الهاتف، وسير معالجة الـ 72 ساعة عند التعرض للاحتيال، والفحص في الوضع الآمن عند الإصابة بالفيروسات، وست خطوات لاستعادة الحسابات المخترقة، والتعامل مع رسائل الفدية، وخمس عادات وقائية يومية.
\end{itemize}

ثانيًا: إضافة فصلين جديدين

\begin{itemize}
  \item الفصل 27: التلاؤم مع كبار السن والرابط الأُسري — جدول مقارنة الوضع المبسط بين العلامات التجارية، وست حِيَل احتيالية وثلاثة خطوط دفاع، والمساعدة عن بُعد (ChangLian/Sunflower)، ووضع الرعاية في WeChat والإدخال الصوتي، وSOS / اكتشاف السقوط / تذكير الأدوية.
  \item الفصل 28: دليل الشراء — طريقة الاختيار في ثلاث خطوات، وجدول مواصفات «التركيز على خمسة بنود فقط»، وتوصيات لخمس فئات ميزانية، وجدول اختيار الهاتف لسبعة سيناريوهات، وخمس فخاخ عند الشراء، وتحليل توقيت الشراء والقنوات.
\end{itemize}

ثالثًا: إضافة ميزات القراءة الصوتية والتفاعل

\begin{itemize}
  \item شريط تحكم عائم للقراءة الصوتية: يدعم التشغيل/الإيقاف المؤقت، والفقرة السابقة/التالية، وضبط السرعة من 0.5x إلى 2.5x، والإيقاف.
  \item زر سماعة سريع على يمين الشاشة: تشغيل القراءة الصوتية بضغطة واحدة، مع البدء تلقائيًا من موضع القراءة الحالي.
  \item اختصارات لوحة المفاتيح: R لتبديل القراءة الصوتية، وESC للإيقاف.
  \item تحسين تجربة القراءة الصوتية: تقدم تلقائي مستمر بين الفصول، وإخفاء عناصر التحكم بالقراءة الصوتية عند الطباعة.
\end{itemize}

رابعًا: تحديث هيكل الكتاب

تحديث متزامن كامل لجدول المحتويات، ومسار التعلم (من أربع خطوات إلى خمس)، وخريطة الفصول، ورقم الإصدار (من 1.6 إلى 1.7)، وترقيم الملاحق (من خمسة إلى ستة)، والبيانات الوصفية.

خامسًا: تحديث حداثة المحتوى

\begin{itemize}
  \item معيار المعالجات: إدراج Snapdragon 8 Elite Gen5 / Dimensity 9500 ضمن الفئة العليا (معيار عام 2026).
  \item مراجع الهواتف الرائدة: iPhone 17 Pro وSamsung S26 Ultra وXiaomi 17 Ultra وHuawei Pura 80 Ultra وvivo X300 Pro.
  \item فصول الذكاء الاصطناعي: التحديث إلى DeepSeek R1 وChatGPT-5 ومعيار قوة NPU على الجهاز بواقع 80 TOPS.
  \item الهواتف القابلة للطي: هبوط نطاق الأسعار إلى فئة 4000-5000 يوان.
  \item دليل الشراء: تحديث كل فئات الميزانية والتوصيات حسب السيناريوهات إلى سوق 2026.
  \item الشبكة والمنظومة: تحديث Wi-Fi 7 / 5G-A / منظومة HarmonyOS / LE Audio كلها إلى أحدث المواصفات.
\end{itemize}

سادسًا: وحدة ميزات جديدة: لوحة مؤشرات التعلم (أسفل منطقة Hero)

\begin{itemize}
  \item مخطط دائري لتقدم القراءة: يعرض النسبة المئوية للفصول المقروءة في الوقت الفعلي، مع حركة SVG دائرية.
  \item فحص ذاتي لصحة الهاتف: قائمة من 8 بنود تحقق أمني (كلمة المرور، والتحقق بخطوتين، والعثور على الهاتف، والنسخ السحابي الاحتياطي، ونسخ WeChat الاحتياطي، وأمان الدفع، وتطبيق مكافحة الاحتيال، وتنظيف الدفع بدون كلمة مرور).
  \item حفظ الحالة: تُكتب حالة التحديد في localStorage، وتُحسب درجة الصحة ديناميكيًا وتُعرض بشريط تقدم متدرج الألوان.
  \item رسائل تنبيه متدرجة: تعرض تغذية راجعة مختلفة حسب النتيجة مثل «الدرجة الكاملة / جيد / مقبول / مخاطر أمنية».
\end{itemize}

\booksection{أبرز مستجدات الإصدار 1.8}

إضافة القسم السابع «برمجة الذكاء الاصطناعي بالتطبيق العملي» (الفصول 37-47، 11 فصلًا): ثورة النموذج وهندسة الوكلاء Agentic Engineering، والمبادئ الخمسة للتفكير في برمجة الذكاء الاصطناعي، وهندسة السياق والتطوير الموجّه بالمواصفات (SDD)، وهندسة المتطلبات، ومشهد أدوات برمجة الذكاء الاصطناعي في 2026، ونظرة شاملة على النماذج الكبيرة والتطبيق العملي لواجهات API، وهندسة الحوار، والدورة الكاملة لمشروع من 0 إلى 1، وضمان الجودة، وطريق الترقي (MCP/تعاون الوكلاء المتعددين)، والدليل المرجعي السريع. مع تحديث متزامن لجدول المحتويات، ومسار التعلم (من خمس خطوات إلى ست)، وخريطة الفصول، وفهرس المراجع السريع. المحتوى مستمد من «عقلية برمجة الذكاء الاصطناعي: دليل منهجي من الصفر إلى التطبيق الشامل» نسخة أغسطس 2026.

إضافة القسم الثامن «تطوير Android بالتطبيق العملي» (الفصول 48-58، 11 فصلًا) والقسم التاسع «تطوير iOS بالتطبيق العملي» (الفصول 59-69، 11 فصلًا): تغطي المنصتان معًا تهيئة البيئة (Android Studio / Xcode)، والتعلم السريع للغة (Kotlin / Swift)، وأطر الواجهات (Jetpack Compose / SwiftUI)، والبيانات والشبكات، ومشروعين عمليين متكاملين (دفتر حسابات + استعلام عن الطقس)، والاختبار والنشر (Google Play / App Store)، والدليل المرجعي السريع. مع تحديث متزامن لجدول المحتويات، ومسار التعلم (من ست خطوات إلى ثمان)، وخريطة الفصول، وفهرس المراجع السريع. المحتوى مستمد من «تطوير Android من الصفر» و«تطوير iOS من الصفر» نسخة أغسطس 2026.

\booksection{إضافات المحتوى: البرمجيات العملية وحجب المزعجات}

\begin{itemize}
  \item يضيف الفصل 13 قسمًا سادسًا جديدًا «أدوات النظام العملية»: يجمع أدوات التنظيف مثل QingZhuo، والتنظيف المدمج في النظام، وتطبيق مركز مكافحة الاحتيال الوطني، ومدير الهاتف من الشركات المصنّعة، وحجب المزعجات من شركات الاتصالات وAdGuard، مع مبدأ اختيار «النظام أولًا / الإضافات الرسمية / الاعتدال مع أدوات الطرف الثالث».
  \item جدول مقارنة حجب المزعجات المدمج في الهواتف حسب العلامة التجارية: يغطي iOS وXiaomi/Redmi وHuawei/Honor وOPPO/realme/OnePlus وvivo/iQOO وSamsung من حيث مسارات الدخول والخيارات الموصى بتفعيلها.
  \item توسيع قدرات حجب المزعجات في الفصل 8: إضافة شرح عملي للتعرف على المتصلين، والقائمتين السوداء والبيضاء، وفلترة الرسائل النصية بالكلمات المفتاحية، وكتم المتصلين غير المعروفين في iOS، والربط بوضع عدم الإزعاج وغير ذلك، مع الربط المتبادل مع الفصل 13.
  \item توصية الحماية ثلاثية الطبقات: حجب النظام + حجب المزعجات من شركة الاتصالات + مركز مكافحة الاحتيال الوطني، مع التأكيد على القائمة البيضاء لتفادي الحجب الخاطئ.
\end{itemize}

\booksection{إضافات المحتوى: تحذيرات الطقس / تحذيرات الكوارث الطبيعية / تحذيرات الزلازل}

\begin{itemize}
  \item يضيف الفصل 13 قسمًا سابعًا: تحذيرات الطقس (القنوات الرسمية للأرصاد + تطبيق الطقس في النظام + مستويات الألوان)، وتحذيرات الكوارث الطبيعية (قنوات متعددة من الأرصاد وإدارة الطوارئ)، ومبدأ عمل تحذيرات الزلازل والمفاهيم الخاطئة الشائعة (التحذير $\neq$ التنبؤ).
  \item جدول تفعيل تحذيرات الزلازل حسب العلامة التجارية: مداخل التفعيل الشائعة ونقاط الأذونات المهمة في Xiaomi وHuawei/Honor وأجهزة OPPO وvivo وiPhone وSamsung وغيرها.
  \item يضيف الفصل 32 «السيناريو صفر»: ترتيب التصرف الآمن بعد إنذار الزلزال، والتعامل مع الأمطار الغزيرة والأعاصير، والربط مع SOS عبر الأقمار الصناعية والاتصال دون شبكة؛ مع ربط متبادل بين دليل الطوارئ وإعدادات التحذير في الفصل 13.
  \item الربط مع قسم الطقس: إضافة إلى بند تطبيق الطقس في قسم السفر إشعارات التحذير وتذكير «لن تصلك التنبيهات إلا بعد تفعيل الإشعارات».
\end{itemize}

\booksection{إضافات المحتوى: التعاون عبر الأجهزة وأدوات الإنتاجية}

\begin{itemize}
  \item القسم الثامن من الفصل 13: التحكم في الكمبيوتر من الهاتف (Sunflower وToDesk) وخطوط الحماية من الاحتيال؛ الذكاء الاصطناعي في العمل (Trae / WorkBuddy / Dumate) وسير عمل التحكم في الكمبيوتر من الهاتف؛ وخلفية «السطح الاصطناعي» المتحركة من miHoYo.
  \item النقل السريع ومشاركة الحافظة عبر الشبكة المحلية: LocalSend وFeiGe وAirDrop؛ وKDE Connect والحافظة الموحدة، والنقل المتبادل بين هواتف وأجهزة الكمبيوتر من نفس العلامة التجارية.
  \item كلمات المرور الديناميكية 2FA: Aegis وMicrosoft Authenticator، مع رموز الطوارئ الاحتياطية ونقاط مهمة عند تغيير الهاتف.
  \item تطبيق TickTick والمسح الضوئي للمستندات الورقية: تفصيل استخدامات الجدول/المهام؛ واستكمال أدوات المسح وحيل التصوير.
\end{itemize}

نتمنى لك أن تتقن هاتفك، وأن تتمتع بسهولة التقنية ومتعتها، وأن تصبح مواطنًا رقميًا واعيًا وواثقًا.

— مرشدك في الحياة الرقمية
